\section*{Capítulo \ref{ch:b3:membrane-traffic} --- Tráfego de membranas e endereçamento de proteínas}

\begin{solution}{exo:b3:membrane-traffic:1}
Trecho hidrofóbico amino-terminal: \omterm{def:b3:membrane-traffic:signals}{peptídeo-sinal}, rumo ao retículo
endoplasmático (e, por padrão, adiante até a superfície). PKKKRKV:
\omterm{def:b3:membrane-traffic:signals}{sinal de localização nuclear}, rumo ao núcleo pelo poro. Hélice
anfipática amino-terminal: \omterm{def:b3:membrane-traffic:signals}{pré-sequência} mitocondrial, rumo à matriz
mitocondrial. KDEL: recuperação do Golgi de volta ao retículo.
Manose-6-fosfato: do Golgi trans rumo ao \omterm{def:b3:membrane-traffic:endocytosis}{endossomo} e ao \omterm{def:b3:membrane-traffic:endocytosis}{lisossomo}.
\end{solution}

\begin{solution}{exo:b3:membrane-traffic:2}
(1) Sem membranas, a cadeia traduzida era cerca de vinte resíduos mais
longa que a proteína secretada: um segmento é removido durante a
secreção. (2) Com membranas presentes durante a tradução, o produto
tinha o tamanho maduro e ficava protegido da protease: ele entrara nas
vesículas e perdera o segmento dentro delas. (3) Membranas
acrescentadas depois da tradução não faziam nem uma coisa nem outra: a
entrada é acoplada à síntese --- é cotraducional --- e o segmento, o
\omterm{def:b3:membrane-traffic:signals}{peptídeo-sinal}, só funciona numa cadeia nascente.
\end{solution}

\begin{solution}{exo:b3:membrane-traffic:3}
Do retículo ao Golgi: \omterm{def:b3:membrane-traffic:coats}{COPII}, para a frente. Do Golgi ao retículo:
\omterm{def:b3:membrane-traffic:coats}{COPI}, retrógrado. Do Golgi trans ao \omterm{def:b3:membrane-traffic:endocytosis}{endossomo}: \omterm{def:b3:membrane-traffic:coats}{clatrina} com seus
adaptadores, para a frente rumo ao \omterm{def:b3:membrane-traffic:endocytosis}{lisossomo}. Da membrana plasmática ao
\omterm{def:b3:membrane-traffic:endocytosis}{endossomo}: \omterm{def:b3:membrane-traffic:coats}{clatrina}, para dentro (endocitose).
\end{solution}

\begin{solution}{exo:b3:membrane-traffic:4}
Os açúcares só são acrescentados dentro do lúmen do retículo e do
Golgi. A face lumenal de cada vesícula continua lumenal ao longo do
brotamento e da fusão, e, quando uma vesícula se funde com a membrana
plasmática, sua face lumenal se torna o exterior da célula; a
orientação nunca se inverte, de modo que o que foi glicosilado por
dentro acaba por fora.
\end{solution}

\begin{solution}{exo:b3:membrane-traffic:5}
$t^{*} = \ln(0.05/0.2)/(0.05 - 0.2) = \ln 4/0.15 = \qty{9.2}{min}$;
$B(t^{*}) = (0.2/(-0.15))(e^{-1.85} - e^{-0.46}) = 0.63$. Aos
\qty{60}{min}, $A \approx 0$, $B = 1.33\,(e^{-3} - e^{-12}) = 0.066$,
logo $C = 0.93$.
\end{solution}

\begin{solution}{exo:b3:membrane-traffic:6}
O \omterm{def:b3:membrane-traffic:signals}{peptídeo-sinal} põe a extremidade amino no lúmen; cada trecho
hidrofóbico é alternadamente uma sequência de parada ou de início de
transferência, de modo que há quatro hélices transmembrana (cerca dos
resíduos 30--52, 90--112, 150--172, 210--232). Segmentos: extremidade
amino lumenal; 52--90 citosólico; 112--150 lumenal; 172--210
citosólico; extremidade carboxi lumenal. Glicosilação só nos segmentos
lumenais --- a extremidade amino, 112--150 e a extremidade carboxi ---,
que se tornam extracelulares na superfície.
\end{solution}

\begin{solution}{exo:b3:membrane-traffic:7}
$J_{\max} = \num{30000}\times 0.25\times 0.05/0.30 = 1250$ por minuto;
fração na superfície $k_{r}/(k_{i} + k_{r}) = 0.05/0.30 = 0.17$. Dobrar
$R$ dobra $J$ exatamente; a taxa de retorno é o gargalo (a maioria dos
receptores está dentro), e dobrar $k_{r}$ daria $\num{30000}\times 0.25
\times 0.1/0.35 \approx 2140$, um fator $1.7$.
\end{solution}

\begin{solution}{exo:b3:membrane-traffic:8}
(a) Sem receptor: as hidrolases marcadas seguem a rota padrão e são
secretadas; os \omterm{def:b3:membrane-traffic:endocytosis}{lisossomos} ficam sem enzimas e se enchem de material não
digerido --- uma doença de depósito. (b) Sem fosfotransferase: o mesmo
fenótipo por outra lesão (doença das células~I), com as enzimas sem
marca e secretadas. (c) \omterm{def:b3:membrane-traffic:endocytosis}{Endossomos} neutralizados: o receptor não
consegue soltar sua carga, e por isso não é reciclado e as hidrolases
são mal entregues ou secretadas; a LDL e outros receptores também param
de ciclar.
\end{solution}

\begin{solution}{exo:b3:membrane-traffic:9}
A tirosina pertence ao \omterm{def:b3:bioinformatics:motif}{motivo} da cauda que o adaptador de \omterm{def:b3:membrane-traffic:coats}{clatrina}
reconhece; sem ela o receptor não é recolhido para as fossetas
revestidas. Os receptores ligam a LDL normalmente, mas ficam espalhados
de modo uniforme pela superfície, excluídos das fossetas, e o ligante
permanece do lado de fora --- uma doença de endereçamento, e não de
ligação.
\end{solution}

\begin{solution}{exo:b3:membrane-traffic:10}
Com $k_{1} = k_{2} = k$, tente $B = kt\,e^{-kt}$:
$\mathrm{d}B/\mathrm{d}t
= k e^{-kt} - k^{2}t e^{-kt} = kA - kB$, como
se pedia, com $B(0) = 0$. O pico está onde $1 - kt = 0$: $t^{*} = 1/k$,
$B = e^{-1} = 0.37$. Em geral,
$C(t) = 1 - (k_{2}e^{-k_{1}t} - k_{1}e^{-k_{2}t})/(k_{2} - k_{1})$, que
é simétrica na troca de $k_{1}$ por $k_{2}$: a curva dos grânulos
sozinha não pode dizer qual dos dois compartimentos é o rápido.
\end{solution}

\begin{solution}{exo:b3:membrane-traffic:11}
Cada vesícula tem área $4\pi(50\,\text{nm})^{2} = \qty{0.031}{\micro
m^2}$; mil por minuto dão \qty{31}{\micro m^2/min}, e portanto toda a
superfície em $3000/31 \approx \qty{95}{min}$. A célula não encolhe
porque a membrana volta por exocitose (vesículas de reciclagem) à mesma
taxa --- um equilíbrio de fluxos. Cada vesícula contém
$5\times 10^{-19}$ L; mil por minuto durante uma hora dão
$3\times 10^{-14}$ L, cerca de \qty{1.5}{\%} do volume de uma célula de
\qty{2}{pL} por hora.
\end{solution}

\begin{solution}{exo:b3:membrane-traffic:12}
A toxina botulínica age no terminal motor: nenhuma acetilcolina é
liberada, o músculo não recebe comando e fica frouxo --- paralisia
flácida. A toxina tetânica é levada pelo axônio motor acima e passa aos
interneurônios inibitórios que normalmente contêm os neurônios motores;
com a \omterm{def:b3:membrane-traffic:fusion}{SNARE} deles clivada, nenhuma glicina ou GABA é liberada, os
neurônios motores disparam sem oposição e todo músculo se contrai ---
paralisia espástica. Alguns nanogramas de toxina botulínica injetados
num músculo hiperativo o silenciam localmente por meses (até que nova
\omterm{def:b3:membrane-traffic:fusion}{SNARE} seja feita e o terminal se recupere): um tratamento para
distonias, espasticidade, estrabismo e rugas.
\end{solution}

\begin{solution}{pb:b3:membrane-traffic:1}
\textbf{1.} $10^{7}\times \qty{25000}{Da}\times 1.66\times 10^{-24}$ g
$= \qty{0.42}{pg}$ por minuto, \qty{600}{pg} por dia --- o dobro de seu
próprio conteúdo proteico.
\textbf{2.} Volume da vesícula $\tfrac{4}{3}\pi\,35^{3} = 1.8\times 10^{5}$
nm$^{3}$, \qty{30}{\%} disso $5.4\times 10^{4}$; uma molécula de enzima
$\tfrac{4}{3}\pi\,2^{3} = 34$ nm$^{3}$: cerca de \num{1600} moléculas
por vesícula; $10^{7}/1600 \approx 6200$ vesículas por minuto.
\textbf{3.} Área $4\pi\,35^{2} = \qty{0.0154}{\micro m^2}$ cada,
\qty{95}{\micro m^2} por minuto: os \qty{6000}{\micro m^2} do retículo
em cerca de uma hora, se nada voltasse.
\textbf{4.} $\qty{95}{\micro m^2} = 9.5\times 10^{7}$ nm$^{2}$, duas
monocamadas a \qty{0.7}{nm^2}: $2.7\times 10^{8}$ lipídios por minuto. A
célula devolve membrana por vesículas \omterm{def:b3:membrane-traffic:coats}{COPI} e por reciclagem, e
sintetiza lipídio para repor o que sai de vez nos grânulos.
\textbf{5.} Volume do grânulo $5.2\times 10^{8}$ nm$^{3}$, \qty{30}{\%}
disso $1.6\times 10^{8}$: $4.7\times 10^{6}$ moléculas por grânulo;
$6\times
10^{8}$ moléculas por hora enchem cerca de $130$ grânulos.
\textbf{6.} Cada grânulo acrescenta $\pi d^{2} = \qty{3.1}{\micro m^2}$;
$300$ acrescentam \qty{940}{\micro m^2} a \qty{50}{\micro m^2}: vinte
vezes mais. A membrana precisa ser recuperada por endocitose
compensatória em minutos.
\textbf{7.} $t^{*} = \ln 2/0.05 = \qty{13.9}{min}$, $B = 0.50$.
\textbf{8.} $A = e^{-3} = 0.05$; $B = 2(e^{-1.5} - e^{-3}) = 0.35$;
$C = 0.60$.
\textbf{9.} \qty{10}{min} e \qty{20}{min}; total médio \qty{30}{min}.
\textbf{10.} $t^{*} = \ln(0.1)/(-0.09) = \qty{25.6}{min}$; $B(t^{*}) =
(0.1/(-0.09))(e^{-2.56} - e^{-0.256}) = 0.77$: o Golgi incha,
abarrotado com três quartos da carga.
\textbf{11.} No pico, $k_{1}e^{-k_{1}t^{*}} = k_{2}e^{-k_{2}t^{*}}$;
substituindo em $B$ obtém-se $B_{\max} = (k_{1}/k_{2})^{\,k_{2}/(k_{2} -
k_{1})}$: a potência é $k_{2}/(k_{2} - k_{1})$. Para $k_{1} > k_{2}$ o
expoente é negativo e a base é maior que $1$; para $k_{1} < k_{2}$ a
base é menor que $1$ e o expoente é positivo --- nos dois casos o valor
fica abaixo de $1$ (confira: $2^{-1} = 0.5$ aqui).
\textbf{12.} Para $k_{1} \gg k_{2}$, $B(t) \to e^{-k_{2}t}$: a marcação
está no Golgi de imediato e sai a $k_{2}$ --- a etapa do retículo fica
invisível e a curva do Golgi é um simples decaimento.
\textbf{13.} $\theta = 2/(2+2) = 0.5$; $J = \num{50000}\times 0.2\times
0.1\times 0.5/(0.1 + 0.1) = 2500$ partículas por minuto, \num{150000}
por hora.
\textbf{14.} $1.5\times 10^{5}\times 1500 = 2.3\times 10^{8}$ moléculas
de colesterol por hora; $5\times 10^{9}/2.3\times 10^{8} \approx 22$
horas.
\textbf{15.} Fração na superfície $k_{r}/(k_{r} + k_{i}\theta) = 0.5$;
uma viagem de ida e volta leva $1/(k_{i}\theta) + 1/k_{r} = 10 + 10 = \qty{20}{min}$: cerca de $60$ viagens em vinte horas.
\textbf{16.} $J = 1250$ por minuto. Abaixo da saturação
$J \propto R\,[L]$, de modo que, com metade dos receptores, a LDL
plasmática precisa dobrar para a mesma depuração.
\textbf{17.} Com $R$ de volta a \num{50000}: a depuração se restabelece
e a LDL plasmática volta rumo ao normal --- o efeito da estatina.
\textbf{18.} Sem receptores, a LDL só é depurada por rotas
inespecíficas lentas, e por isso circula por dias em vez de horas, é
oxidada e é captada pelos receptores depuradores dos macrófagos, que se
tornam as células espumosas das placas ateroscleróticas; e as
estatinas, que funcionam induzindo receptores, não têm o que induzir.
\textbf{19.} $V = \tfrac{4}{3}\pi(0.25)^{3} = \qty{0.065}{\micro m^3} =
6.5\times 10^{-17}$ L. A pH $7.2$: $6.3\times 10^{-8}\times 6.5\times
10^{-17} = 4\times 10^{-24}$ mol --- cerca de $2$ prótons livres. A pH
$4.7$: $2\times 10^{-5}\times 6.5\times 10^{-17} = 1.3\times 10^{-21}$
mol, cerca de $780$.
\textbf{20.} $50\times 780 \approx \num{39000}$ prótons; $50$ bombas a
$200$ por segundo movem \num{10000} por segundo: cerca de \qty{4}{s}.
\textbf{21.} Bombear carga positiva para dentro deixa o lúmen positivo,
e alguns milhares de cargas parariam a bomba; o cloreto entra por um
canal (ou cátions saem) para neutralizar a carga.
\textbf{22.} Enzimas inativas em pH neutro não fazem mal se um \omterm{def:b3:membrane-traffic:endocytosis}{lisossomo}
vaza ou se elas são mal endereçadas ao citosol ou à superfície: a
exigência de acidez confina a digestão ao compartimento construído para
ela.
\textbf{23.} As duas precisam carregar um sensor de pH --- histidinas,
cuja protonação perto de pH~6 muda a superfície de ligação: o receptor
de LDL dobra um domínio sobre seu próprio sítio de ligação quando
protonado, e o receptor de manose-6-fosfato perde afinidade do mesmo
modo.
\textbf{24.} $10^{7.2 - 4.7} = 10^{2.5} \approx 300$ vezes de
concentração; a base aprisionada consome prótons, eleva o pH lisossomal
e inativa as hidrolases --- que é como a cloroquina mata os parasitos da
malária em seu vacúolo alimentar e como ela bloqueia a \omterm{def:b3:membrane-traffic:endocytosis}{autofagia} nos
experimentos.
\textbf{25.} Pico do Golgi aos \qty{14}{min}; umas \num{6200} vesículas
deixam o retículo por minuto; \num{150000} partículas de LDL captadas
por hora.
\end{solution}
